% Generated draft: requires mathematical and source review.
\documentclass[11pt]{article}
\usepackage[margin=25mm]{geometry}
\usepackage{fontspec}
\setmainfont{Latin Modern Roman}
\usepackage{amsmath,amssymb,mathtools,unicode-math}
\setmathfont{Latin Modern Math}
\usepackage{enumitem,needspace,xurl,hyperref}
\setlength{\parindent}{0pt}
\setlength{\parskip}{5pt}
\setlength{\emergencystretch}{3em}
\newcommand{\N}{\mathbb{N}}
\newcommand{\Z}{\mathbb{Z}}
\newcommand{\Q}{\mathbb{Q}}
\newcommand{\R}{\mathbb{R}}
\newcommand{\C}{\mathbb{C}}
\newcommand{\F}{\mathbb{F}}
\newcommand{\A}{\mathbb{A}}
\newcommand{\PP}{\mathbb{P}}
\newcommand{\E}{\mathbb{E}}
\newcommand{\eps}{\varepsilon}
\newcommand{\dd}{\,\mathrm{d}}
\begin{document}
\section*{UnsolvedMath ID: 8800091}
\textbf{Source problem code:} catalogue entry
\subsection{Definitions and mathematical statement}
In the usual imaginary-quadratic setting, let $K$ be an imaginary quadratic field with maximal order $\mathcal O_K$. Its order of conductor $f$ is
\[
 \mathcal O_f=\Z+f\mathcal O_K.
\]
It is nonmaximal when $f>1$. CM points associated with elliptic curves whose endomorphism ring is $\mathcal O_f$ give the standard source of points from nonmaximal orders; calling particular images Heegner points requires the relevant modular curve, Heegner hypotheses, and any modular parametrization.

The problem asks to prove analogues of previously stated Heegner-point results for these nonmaximal-order points. Those antecedent results are absent from the supplied record. Their conclusion, field and level hypotheses, distribution or independence statement, constants, and parameters being varied cannot be recovered.

Consequently the exact retained task is extension of the omitted source results from their intended maximal-order setting to nonmaximal orders. It is not a claim that nonmaximal-order CM or Heegner points have no construction, nor a guessed distribution theorem. A precise formulation needs the original theorem and its conductor, discriminant, level, and field conventions; none is silently selected.
\subsection{Short English statement}
Extend the Heegner-point theorems cited in the workshop to points associated with nonmaximal orders. The orders can be described by their conductor, but the actual results to be generalized are missing.

The supplied summary reports no verified matching analogue and speculates about distributional bounds. That assessment does not identify the absent theorem or establish a general lack of nonmaximal-order theory. The OPEN-TRIAGE label and 2006 source are retained with current primary-literature audit required.
\subsection{Sources}
\begin{itemize}
\item Local UnsolvedMath registry, permanent ID 8800091, source code catalogue entry, ``Heegner points from nonmaximal orders,'' proposed by Michael Rosen and Joseph Silverman. Dataset curation and metadata: Ulam AI and UnsolvedMath Contributors, CC BY 4.0.
\item IPAM, \emph{Workshop on Number Theory and Cryptography --- Open Problems} (2006), Attachment 6669, PDF page 22, printed slide 16, schedule locator \url{https://www.ipam.ucla.edu/programs/workshops/workshop-i-number-theory-and-cryptography-open-problems/?tab=schedule}; pdf-and-legacy-ppt visual-audit extraction, accessed 2026-07-29.
\item The referenced antecedent Heegner-point results and their bibliographic entries are absent; no exact later resolution is supplied.
\item Source provenance reports inspection of 17 PDFs and two legacy PPT decks among 19 surviving attachments, excluding repetitions and solved definitions; this is supplied provenance, not a new inspection claimed here.
\item Current primary-literature status audit and redistribution terms NEEDS\_REVIEW for the publicly accessible author/publisher copy, separately from the dataset licence.
\end{itemize}
\end{document}
